% 图片致谢。本卷所用照片的图片来源说明（第 1 册）。
% 机器可读的许可记录位于 images/book1/CREDITS.md；增加或删除照片时，请保持两个文件同步。
\chapter*{图片致谢}

本书中的插图为原创。照片依其各自的自由许可使用，谨向其作者致谢：

\begin{itemize}
  \item \emph{约翰·道尔顿}（原子与分子一章）：Charles Turner 依 Joseph Allen 之作所刻版画，1834 年，美国国会图书馆经维基共享资源提供，公有领域。
  \item \emph{道尔顿 1808 年元素表}（原子与分子一章）：摘自约翰·道尔顿《化学哲学新体系》的插页，由互联网档案馆扫描并经维基共享资源提供，公有领域。
  \item \emph{铝土矿}（原材料一章）：Werner Schellmann 摄，经维基共享资源提供，CC BY-SA 2.5。
  \item \emph{花岗岩}（纯物质一章）：Eurico Zimbres 摄，经维基共享资源提供，CC BY-SA 2.0 br。
  \item \emph{无水硫酸铜}（物质的鉴别一章）：W.~Oelen 摄，经维基共享资源提供，CC BY-SA 3.0。
  \item \emph{硫酸铜晶体}（物质的鉴别一章）：Stephanb 摄，经维基共享资源提供，CC BY-SA 3.0。
  \item \emph{安托万-洛朗·拉瓦锡与玛丽-安妮·帕尔泽·拉瓦锡}（配平方程式一章）：Jacques-Louis David 绘于 1788 年，大都会艺术博物馆经维基共享资源提供，公有领域。
  \item \emph{贝克莱特电话机}（塑料一章）：Holger Ellgaard 摄，经维基共享资源提供，CC BY-SA 3.0。
  \item \emph{欧内斯特·卢瑟福}（原子内部一章）：照片，George Grantham Bain 收藏，美国国会图书馆，经维基共享资源提供，公有领域。
  \item \emph{岩盐}（离子一章）：Robert M.~Lavinsky 摄，经维基共享资源提供，CC BY-SA 3.0。
  \item \emph{自由女神像}（金属与腐蚀一章）：Elcobbola 摄，经维基共享资源提供，公有领域。
  \item \emph{德米特里·门捷列夫}及其\emph{1869 年的元素表}（元素周期表一章）：经维基共享资源提供，公有领域。
  \item \emph{蟹状星云}（元素一章）：NASA、ESA、J.~Hester 与 A.~Loll，公有领域。
  \item \emph{钠}（电子壳层一章）：Dnn87 摄，经维基共享资源提供，CC BY-SA 3.0。
  \item \emph{安瓿瓶中的氯}（电子壳层一章）：W.~Oelen 摄，经维基共享资源提供，CC BY-SA 3.0。
  \item \emph{阿梅迪奥·阿伏伽德罗}（摩尔一章）：石版画，依 C.~Sentier 1856 年的素描绘制，Edgar Fahs Smith 收藏，经维基共享资源提供，公有领域。
  \item \emph{雪晶}（极性一章）：Wilson Bentley 摄于 1902 年，经维基共享资源提供，公有领域。
  \item \emph{路易·巴斯德}（立体化学一章）：Paul Nadar 摄，经维基共享资源提供，公有领域。
  \item \emph{酒石酸铵晶体}（立体化学一章）：Marie-Lan Taÿ Pamart 摄，经维基共享资源提供，CC BY 4.0。
  \item \emph{斯万特·阿伦尼乌斯}（酸与碱一章）：1909 年的凹版照相版，经维基共享资源提供，公有领域。
  \item \emph{伏打电堆}（电池一章）：GuidoB 摄，经维基共享资源提供，CC BY-SA 3.0。
  \item \emph{华莱士·卡罗瑟斯在实验室}（聚合物一章）：摄影者不详，经维基共享资源提供，公有领域。
\end{itemize}

本书仓库中的 \texttt{images/book1/CREDITS.md} 列出了每张照片完整的来源链接与许可网址。
危险象形图采用全球化学品统一分类和标签制度（GHS）的图样，即 \TeX\ Live 中 \texttt{ghsystem} 宏包所绘的版本。

\bigskip
除照片与原创插图外，部分图表为 AI 生成的插图，使用 OpenAI 图像生成功能、依据本书编辑所撰写的提示词制作，并在收录前经过化学准确性审核。每张生成图像的完整提示词均记录于仓库中的 \texttt{images/book1/ai/PROMPTS.md}。
\clearpage
